\section{Map the Team to the Field}
\label{sec:map}

Questions three and four of Newell ask where a team can work, and
which important issues live there.
For ``where,'' we use the nine areas of the ICSE'27 research track
call for
papers:\footnote{conf.researchr.org/track/icse-2027/icse-2027-research-track}
AI for SE; analytics; architecture and design; dependability and
security; evolution; human and social aspects; requirements and
modeling; SE for AI; and testing and analysis.

We score each area on the same nine constructs as the people.
The areas then join Figure~\ref{fig:focus} as reference elements (the
italic columns).
Each area column is an archetype.
It shows how the usual highly cited researcher of that area scores.
We judged the archetypes from the call for papers, not from a corpus.
\note{claude}{upgrade path: make each area column from its ICSE
distinguished papers of the last five years, coded like the author
papers; then people and areas have one source}

The red tree then answers a question that a topic list cannot answer:
who stands nearest to which area.
Some pairs are strong.
Menzies stands at 92\% match to analytics.
Di Nucci stands at 89\% to evolution.
Schmid stands at 89\% to requirements and to architecture, and
Esposito stands at 86\% to SE for AI.
Some pairs are loose.
The nearest areas of Lenarduzzi (evolution, architecture) reach only
72\%, and the nearest area of Armenti (human and social aspects)
reaches 69\%.
Thus, those two members are less typecast than the rest of us.
Thus, question three answers itself: the natural territory of the team
is the set of areas adjacent to a member.

The far columns are also important.
No author is nearer than 85\% to security (best: Esposito, 78\%) or to
testing (best: 72\%).
Those two areas hold to each other, not to a person.
That is the priced blind spot of this team.
Papers there would go against the grain, or the team needs a new
collaborator who brings that column nearer.

The red tree also shows a larger shape: the team divides in two.
One wing measures.
That wing is Menzies and Esposito, together with analytics and the two
AI areas.
One wing crafts.
That wing is Schmid, Armenti, Di Nucci, and Lenarduzzi, together with
architecture, requirements, human aspects, and evolution.
This split is also a division of labor.
When a paper observes a corpus and builds a thing, each wing takes one
half.

Last, the grid can propose a topic for the full team.
Three calculations use the same matrix: (a)~the area nearest to the
centroid of the six persons (SE for AI and evolution, tied at 83\%);
(b)~the area with the best worst member (evolution, worst member at
64\%); (c)~the midpoint of each pair of areas, scored the same way,
where the best pairs (analytics plus human aspects at 68\%; AI for SE
plus human aspects, average 76\%) beat each single area.
Thus, blends win.
A topic that crosses the two wings serves this team better than one
area.
The best blend is AI for SE crossed with human and social aspects.
That blend describes the paper that you now read.

\begin{pause}{the topic}
Ask: do the archetype columns agree with your reading of the call for
papers?
Would you spend money on the winning blend?
Or must a new collaborator repair the blind spot (security,
testing)?\\[0.2em] To change the outcome: edit the nine area columns
in {\ttfamily focus.py} (or use the area list of a different venue),
then do {\ttfamily make focus}; the centroid, maximin, and blend
scores follow.
\end{pause}
